A présent, nous allons nous intéresser aux limites liées aux données utilisées dans les visualisations, en particulier la perte ou l'altération d'informations qu'elles peuvent entraîner.

Lorsque nous travaillons sur des documents tels que des manuscrits, il est souvent nécessaire de les transformer par des processus comme la numérisation, l'extraction d'image et l'annotation. Ces actions permettent à la fois de conserver les documents originaux en facilitant leur accès numérique mais aussi de mettre en lumière des connaissances et éléments jusque-là cachés. Par la suite, nous pouvons utiliser ces informations pour réaliser des visualisations. Cependant, cela modifie le contenu, la structure et la forme matérielle du document, ce qui peut nuire à l'accès à sa réelle signification\footcite{vancisinExternalizingTransformationsHistorical2020}. La numérisation en elle-même fait également perdre une partie de la matérialité du document\footcite{maBoundariesExtensionsChallenges2022}.

Dans le cadre de nos visualisations, nous allons principalement nous concentrer sur les résultats des traitements automatiques comme l'extraction et la reconnaissance de similarités. Nous serions alors tentés de négliger toutes les informations relatives au contexte dans lequel ces données s'inscrivent pour ne travailler que sur les images. Cependant, agir de la sorte serait une erreur car nous perdrions des informations essentielles. 

Quand nous entrons un \textit{witness} et sa numérisation dans la base de données, nous perdons de nombreuses métadonnées et informations liées à la matérialité de l'objet.

Lorsque nous lançons un traitement automatique, nous n'obtenons en résultats que des images et des informations qui leur sont associées, comme le score de similarité ou leur lien avec d'autres images. Néanmoins, il est important de ne pas délaisser certaines métadonnées afin de pouvoir continuer à les inscrire dans un contexte. Par exemple, le lieu de conservation du \textit{witness} est une information primordiale pour retracer sa circulation tout comme la date et le lieu de copie. A cela s'ajoute la langue qui permet de distinguer différentes traditions. Il est donc primordial de ne pas délaisser ces informations dans les visualisations que nous allons créer.